\chapter{アルゴリズムの基礎}
\label{ch:seminar-02}

離散 Kantorovich 問題
\[
  \MKD_{\mathbf{C}}(\mathbf{a}, \mathbf{b})
  \defeq \min_{\mathbf{P} \in \CouplingsD(\mathbf{a}, \mathbf{b})}
  \inner{\mathbf{C}}{\mathbf{P}}
\]
は $nm$ 個の変数と $n + m$ 本の等式制約を持つ線形計画であり，
\textbf{ネットワーク単体法}（network simplex）をはじめとする
LP アルゴリズムで厳密に解ける．
ネットワーク単体法は輸送多面体の頂点間を移動しながら
コストを単調に減少させ，多項式時間で最適解に到達する（Orlin, 1997）．

しかし，学習・大規模応用の観点ではいくつかの根本的な制約をもつ：

\begin{itemize}
  \item \textbf{計算量}：$O(nm)$ 個の変数に対する反復・データ構造管理は，
    $n, m$ が $10^5$ オーダーになると現実的でない．
    GPU 並列化にも本質的に向かない（離散探索的）．
  \item \textbf{微分不可能性}：最適値 $\MKD_{\mathbf{C}}(\mathbf{a}, \mathbf{b})$ は
    $\mathbf{a}, \mathbf{b}$ について区分線形にしかならず，
    滑らかな損失関数として自動微分による勾配最適化に組み込みにくい．
  \item \textbf{標本バイアス}：母集団からの有限サンプルに基づいて
    $\MKD_{\mathbf{C}}$ を推定すると，
    高次元では母集団値への収束が指数的に遅い（次元の呪い）．
\end{itemize}

これらの困難は，目的関数に \textbf{エントロピー} を加えて
\textbf{正則化}することで大きく緩和される：
正則化された最適輸送（\textbf{Sinkhorn アルゴリズム}）は
計算が GPU 並列で軽く，最適値は marginals に関して滑らか（微分可能）であり，
さらに統計的安定化の理論的基礎も提供する．
